\documentclass[12pt]{article}
\usepackage[utf8]{inputenc}
\usepackage[T1]{fontenc}
\usepackage{hyperref}
\title{Weyto Language Resources for K-12 Curriculum Development}
\author{Compiled by Claude Sonnet 5.5 (claude-sonnet-5-5) for Andr\'as Kornai}
\date{}
\begin{document}
\maketitle

\section{Introduction}
Weyto (ISO 639-3 woy; Glottocode weyt1237) is described as an extinct or dormant language, said to have been spoken around Lake Tana, Ethiopia, by a small hippopotamus-hunting group. The group now speaks Amharic. Evidence for the language is very thin. The earliest record is from James Bruce, around 1770, who said it was radically different from other languages of Abyssinia but gave no usable data. Eugen Mittwoch (1907) described the people as speaking a dialect of Amharic. The language remains unclassified, though a link to Agaw has been suggested. The Weyto caste still lives along the lake. Searches found no teaching materials, so a curriculum in the language itself is not currently possible.

\section{Primary Educational Resources}
None found. No primer, textbook or school programme for Weyto appeared in search results.

\section{Linguistic and Reference Materials}
\begin{itemize}
\item Wikipedia, ``Weyto language'': \url{https://en.wikipedia.org/wiki/Weyto_language}
\item Wikipedia, ``Weyto caste'' (background on the community): \url{https://en.wikipedia.org/wiki/Weyto_caste}
\item Glottolog entry for Weyto, which lists bibliographic references: \url{https://glottolog.org/resource/languoid/id/weyt1237}
\item Glottolog reference record for Tecle Haimanot, Gebre Selassie (1984), an Addis Ababa University thesis on the Wayto of Lake Tana (only the bibliographic record was seen): \url{https://glottolog.org/resource/reference/id/318616}
\end{itemize}
A 2010 state-of-the-art overview by Chlo\'e Darmon is mentioned in Glottolog, but no URL for it was seen in search results.

\section{Religious and Community Resources}
None found.

\section{Digital and Audio Resources}
None found for the language itself. A Rosetta Project entry lists name variants (Wayto, Weyt'o): \url{https://rosettapanglossia.longnow.org/index.php?title=Woy}

\section{Curriculum Development Context}
No usable language data appear to exist. Any school work would have to be about the Weyto community's history and its Amharic-speaking present, not instruction in Weyto. Further archival research on the 1984 thesis and the works cited in Glottolog would come first.

\end{document}